\subsection{Visão Arquitetural}
\label{sec:arquitetura}

A arquitetura da plataforma de processamento de mídias foi concebida como um \textbf{Monólito Modular Orientado a Eventos}, fundamentada nos princípios do \textit{Domain-Driven Design} (DDD) e \textit{Command Query Responsibility Segregation} (CQRS). Esta abordagem visa maximizar a flexibilidade, escalabilidade e manutenibilidade do sistema, permitindo a evolução independente de suas funcionalidades e a fácil adaptação a novas demandas.

\subsubsection{Monólito Modular e Domain-Driven Design (DDD)}
O sistema é estruturado em \textit{Bounded Contexts} (Contextos Delimitados) claros, como \texttt{auth}, \texttt{media\_management}, \texttt{transcription}, \texttt{summarization}, \texttt{notifications}, \texttt{analytics} e \texttt{metrics}. Cada módulo é auto-contido, com suas próprias camadas de API, Aplicação, Domínio e Infraestrutura, minimizando o acoplamento entre eles e promovendo alta coesão interna. O domínio de cada módulo é rico, encapsulando regras de negócio e garantindo a integridade dos dados.

\subsubsection{Command Query Responsibility Segregation (CQRS)}
O padrão CQRS foi aplicado consistentemente em todos os módulos de negócio para separar as responsabilidades de escrita (comandos) e leitura (consultas).
\begin{itemize}
    \item \textbf{Comandos (Commands)}: Representam intenções de mudança de estado no sistema (e.g., \texttt{CreateMediaCommand}, \texttt{ProcessTranscriptionCommand}). São objetos imutáveis que encapsulam todos os dados necessários para a operação.
    \item \textbf{Manipuladores de Comando (Command Handlers)}: Classes dedicadas (e.g., \texttt{CreateMediaCommandHandler}) que executam a lógica de negócio associada a um comando específico. Eles orquestram as operações de domínio, interagem com repositórios para persistência e publicam eventos de domínio.
    \item \textbf{Consultas (Queries)}: Classes dedicadas (e.g., \texttt{MediaQueries}, \texttt{TranscriptionQueries}) otimizadas para operações de leitura. Elas dependem de repositórios para buscar dados e não modificam o estado do sistema.
    \item \textbf{Fachadas de Serviço (Service Facades)}: As classes de serviço (e.g., \texttt{MediaService}, \texttt{AuthService}) atuam como fachadas finas na camada de aplicação, recebendo requisições da API e despachando os comandos ou eventos apropriados para seus respectivos manipuladores, sem conter lógica de negócio complexa.
\end{itemize}
Esta separação permite otimizar e escalar as operações de leitura e escrita de forma independente, além de simplificar a lógica de cada componente.

\subsubsection{Pipeline de Processamento de Mídia Orientado a Eventos}
O coração da plataforma é um pipeline assíncrono e orientado a eventos, que permite o processamento genérico de diferentes tipos de mídia (vídeo, áudio, texto/PDF).
\begin{itemize}
    \item \textbf{Event Bus}: A comunicação entre os módulos é realizada através de um \textit{Event Bus} (implementado com Redis e Celery), que transporta \textit{Domain Events} tipados (e.g., \texttt{MediaUploaded}, \texttt{AudioExtracted}, \texttt{TextContentReady}).
    \item \textbf{Fluxo de Eventos}:
    \begin{enumerate}
        \item \texttt{media\_management} recebe um upload (\texttt{POST /media}) e publica \texttt{MediaUploaded}.
        \item Se \texttt{MediaUploaded} for do tipo \texttt{VIDEO}, o módulo \texttt{video\_processing} consome o evento, extrai o áudio e publica \texttt{AudioExtracted}.
        \item Se \texttt{MediaUploaded} for do tipo \texttt{AUDIO} ou \texttt{AudioExtracted} for publicado, o módulo \texttt{transcription} consome o evento, transcreve o áudio e publica \texttt{TextContentReady}.
        \item Se \texttt{MediaUploaded} for do tipo \texttt{TEXT}, o módulo \texttt{text\_processing} consome o evento, extrai o texto e publica \texttt{TextContentReady}.
        \item O módulo \texttt{summarization} consome \texttt{TextContentReady} (independentemente da origem), gera um resumo e finaliza o pipeline.
    \end{enumerate}
\end{itemize}
Este design promove um desacoplamento extremo, onde cada módulo reage a eventos genéricos e não precisa conhecer os detalhes de implementação dos módulos adjacentes.

\subsubsection{Padrão State para Gerenciamento do Ciclo de Vida}
A entidade \texttt{Media} implementa o \textit{State Pattern} para gerenciar seu ciclo de vida (e.g., \texttt{UPLOADED}, \texttt{PROCESSING}, \texttt{COMPLETED}, \texttt{FAILED}).
\begin{itemize}
    \item \textbf{Classes de Estado}: Cada estado é representado por uma classe dedicada (e.g., \texttt{UploadedState}, \texttt{ProcessingState}), que encapsula a lógica e as transições permitidas para aquele estado.
    \item \textbf{Delegação}: A entidade \texttt{Media} delega as ações de transição (e.g., \texttt{media.process()}) para seu objeto de estado atual, garantindo que apenas transições válidas ocorram e centralizando a lógica de estado.
\end{itemize}
Este padrão aumenta a robustez do sistema, prevenindo estados inválidos e simplificando a lógica de negócio.

\subsubsection{Seleção Dinâmica de Provedores (Padrão Strategy e Factory)}
Para serviços de infraestrutura (armazenamento, transcrição, sumarização), o sistema utiliza um padrão de \textit{Strategy} e \textit{Factory} para permitir a seleção dinâmica de provedores em tempo de execução.
\begin{itemize}
    \item \textbf{Interfaces Abstratas}: Cada tipo de serviço (e.g., \texttt{StorageService}, \texttt{ISpeechRecognition}) é definido por uma interface abstrata.
    \item \textbf{Registro Dinâmico}: Implementações concretas (e.g., \texttt{S3StorageService}, \texttt{WhisperTranscriber}) se registram automaticamente em factories centrais usando decoradores.
    \item \textbf{Factory Functions}: Funções de fábrica (e.g., \texttt{create\_storage\_service}) criam instâncias de provedores sob demanda, baseadas na escolha do usuário ou em configurações.
\end{itemize}
Isso confere ao sistema alta flexibilidade, permitindo a fácil adição ou substituição de provedores de serviço sem modificar a lógica de negócio principal.

\subsubsection{Raiz de Composição Explícita e Descentralizada}
A injeção de dependência é gerenciada de forma descentralizada, com o arquivo \texttt{main.py} atuando como a única \textbf{Raiz de Composição} da aplicação.
\begin{itemize}
    \item \textbf{Bootstraps por Módulo}: Cada módulo de negócio possui seu próprio arquivo \texttt{bootstrap.py}, responsável por construir e configurar seus serviços, handlers e queries internos.
    \item \textbf{Orquestração Explícita}: O \texttt{main.py} orquestra a inicialização de cada módulo e injeta explicitamente as poucas dependências que cruzam os \textit{Bounded Contexts} (e.g., \texttt{UserRepository} do \texttt{auth} no \texttt{notifications}).
\end{itemize}
Esta abordagem elimina um container centralizado, reduzindo o acoplamento no ponto de configuração e tornando o grafo de dependências da aplicação transparente.

\subsubsection{Observabilidade e Resiliência}
\begin{itemize}
    \item \textbf{Observabilidade Completa}: O sistema integra os três pilares da observabilidade: \textit{logs} estruturados (\texttt{structlog}), \textit{métricas} Prometheus e \textit{tracing} distribuído com OpenTelemetry e Jaeger, permitindo o monitoramento detalhado e a depuração de fluxos assíncronos complexos.
    \item \textbf{Padrão Circuit Breaker}: Chamadas a serviços externos (e.g., provedores de IA) são protegidas pelo \textit{Circuit Breaker Pattern}, prevenindo falhas em cascata e aumentando a resiliência do sistema a indisponibilidades externas.
\end{itemize}
A combinação desses padrões resulta em uma plataforma robusta, escalável e de fácil manutenção, preparada para lidar com as complexidades do processamento de mídia em larga escala.